\section{传输、恢复审计与评价生命周期}
\subsection{所有真实模型调用的唯一准入链}
live 模型经 \code{OpenAICompatibleChatModel.complete} 进入测试门禁，然后生成 request ID、记录 model start、发出请求、记录原始响应与 usage、解析 reply 并记录成功或错误。scripted 模型用于确定性离线机制检查，运行 provenance 由实际模型类型核验，不能通过文字字段伪装为 live。

同一进程首次真实请求前自动运行完整 \cmd{uv run --no-sync pytest}。指纹覆盖 src、tests、scripts、configs、validation 以及依赖声明、AGENTS 和测试契约。当前完整门禁有 900 秒有限截止时间。进程启动后、测试期间或准入后源码指纹变化，都阻止继续请求。

通过状态只保存在该进程内存；本地 gate log 是审计产物，不是可重复使用的 pass stamp。测试失败后该进程不按剩余 trial 自动重复尝试准入。发请求需要实际配置中的实验授权，测试通过本身不授权新增预算、provider fallback 或重跑。
\loc{runtime/model.py; runtime/admission.py}

\subsection{测试隔离与凭据处理}
门禁子进程清除 API key、token、secret、password 等凭据环境项，禁用自动 pytest plugin 加载，并移除会改变测试选择或 provider 的部分环境设置。pytest 使用 scripted 模型、mock HTTP 及项目拥有的 loopback 夹具。

真实模型的 endpoint、model、temperature、超时及输出上限由显式运行配置与提供商组装获得，凭据在运行入口加载。模型响应和传输错误的公开文本与事件对当前配置的 API 密钥值脱敏；源码归档和技术报告不保存凭据。模型 ID 使用实际返回值校核，不能仅凭请求配置名字认定上游实际模型。
\loc{runtime/provider.py; runtime/model.py::_record_response}

\subsection{先记录 usage，再验证响应内容}
模型原始响应中的 usage 在 content、tool JSON 或最终 reply schema 校验前写入日志。因此，格式错误且具有 reported usage 的响应仍计入成本。请求数在 model start 计数；无响应、解析失败和中断都不被解释为零请求。

事件按 request ID 对账，检查重复 start、孤立响应、usage 冲突、缺失 reported usage 和无效响应。未知 usage 与未配置价格不等于数值零；只有所有相关分量已知时，才能给出完整总和。报告使用 provider reported tokens，而文本字节预算单独报告。

HTTP 错误、空 choices、非法工具 JSON 和超时显式失败，不作为 SAFE 票。当前传输不会因这种失败自动换 provider 或执行静默重试；若协议允许的纠正循环再发请求，该请求必须重新进入预算与记账链。

\subsection{耐久事件与全局停止}
Journal 按 JSONL 追加事件，flush 并 fsync 后才继续操作。并发实验使用锁保护日志和预算，单个 trial 保存模型、工具、角色、调用参数、允许工具与准入路径等。结果同时保留原始轨迹和纯汇总，异常不会抹去此前请求。

通用 Budget 约束请求数和 wall-time，LockedBudget 约束共享并发。场景实验另有累计 token 账本、阶段上限、保守 reservation 和硬停止条件。实验取消阻止后续请求进入；线程任务的实际已发生操作仍须入账。

源码快照、配置、provider 参数、gate log、事件、结果、usage、summary 和验收形成可审计产物链。fingerprint 与归档按已授权边界机械读取文件，不能被当作检测语义访问或用来恢复模型可见标签。
\loc{runtime/journal.py; runtime/budget.py; runtime/snapshots.py; evaluation/experiment_guard.py}

\subsection{配置与实验组织}
评价层声明数据角色、实验配置和产物生命周期，与核心运行时分开。
\subsubsection{配置按职责组合}
配置目录按 models、datasets、profiles、experiments、preparation、diagnostics 与 history 分类。dataset 声明样本、角色、来源和前提；profile 声明系统、预算、检索、选择格子和预登记弃权；model 配置声明模型默认参数。专项实验维持自身明确的类型化配置。

advisory 数据与运行 profile 经 \code{compose_advisory_config} 组装成既有类型化协议，字段不能任意交叉覆盖，缺失或无效配置直接报错。解析后的完整配置进入 metadata。文件名中的历史模型名不能替代实际返回模型 ID；含 heldout 字样的旧路径也不能恢复数据盲测资格。

旧配置和历史产物被保留用于版本比较。迁移等价性测试核对明确路径替换与新增字段，不能通过覆盖旧结果来制造一致。准备程序具有修改数据的行为；查看项目架构不授权重建数据或开启实验。
\loc{evaluation/datasets/composition.py; evaluation/datasets/advisory_config.py; runtime/provider.py}

\subsubsection{候选执行与结果状态}
\code{run_candidate} 组装 index、candidate、系统、模型、记录工具和预算，调用 \code{AgenticPipeline.run}。结果为实质判决时记为 completed，为合法弃权时记为 abstained；传输、schema、内部异常或超时有明确 failed/interrupted/not-run 状态。

尚未发送模型请求就被预算阻止，可以记为 not-run；已经发生请求后失败，仍保存请求和工具尝试计数。中断会记录当前结果与时长再传播，不能在汇总中当作没有运行。所有实际执行任务包含在最终模型和工具成本中。
\loc{evaluation/execution.py; evaluation/results.py}

\subsubsection{实验生命周期}
advisory 生命周期先校验配置、模型和数据准备，创建独立 run 目录、源码与配置快照及耐久事件，再按预登记的格子顺序或显式并发执行任务。每格结果增量保留，结束或停止后由纯汇总逻辑输出 usage、summary、acceptance 等文件。

不同 protocol 按各自样本前提检查验收：开发门禁可声明有意弃权；配对检查良性功能保留及预期拒绝；场景实验分别审计 dataset、model 和 developer-scripted 来源。通用测试通过不意味着这些专项前置条件已经满足。

源码发现 runner 的 checkpoint 保存清单身份、源码分析指纹、已完成提交和解析错误。恢复需要相同 detector manifest、数据身份和代码指纹；不兼容状态拒绝恢复。临时 worktree 用于减少同时保存大量完整 checkout 的空间压力，原始候选清单和哈希继续保留。
\loc{evaluation/lifecycle.py; evaluation/runners/run_vulngym_discovery.py; evaluation/repository.py}

\subsubsection{检测输入与评价标签分离}
一般仓库输入可以只包含 repository URL 和 commit；候选定位及性质应由检测流程产生。评价侧 reference 含标签、关键位置和人工核验信息，只有在候选或 findings 已产生后才参与匹配与计分。oracle-diagnostic 使用已知入口作为 seed 时，必须明确其角色，不能报告端到端发现召回。

当前源码发现 evaluator 按规范化仓库、commit、路径和行号精确匹配 critical-operation 位置。旧目标规格中出现的其他行号容忍窗口，不自动成为该实际计分器规则。没有核验负例的正例参考集不能估计 FPR；未匹配的 VULNERABLE finding 也不能仅凭参考表缺项判为假阳性。

开发样本、已暴露的公告路径或 oracle-seeded 数据保持 \code{claim_eligible=false}。架构方披露已查询部分 held-out 评价侧元数据，VulnGym v4 的命令注入相关结论、Python 类别组成和语言组成不再视为盲测。技术报告仅引用已归档披露，不重新读取受保护输入或 evaluator 标签。\cite{disclosure}

根目录 \code{.ignore} 使 ripgrep 默认搜索跳过受保护输入及其源码快照副本；它是默认导航的预防措施，门禁指纹和源码归档继续按独立规则机械读取。搜索排除规则不应被误写成文件本身已经删除或无法被显式访问。

\subsubsection{运行完整性与研究 claim readiness}
\code{assess_run_readiness} 对 acceptance、metadata、summary、usage 及结果/事件中的 fixture 使用作交叉核对。它不只读取一个存储的 passed 布尔。

\begin{table}[htbp]
\centering\small
\begin{tabularx}{\textwidth}{L{37mm}Y}
\toprule 关口 & 实际含义\\\midrule
完整 pytest & 当前代码可进入受门禁约束的实验；未证明检测效果。\\
专项 acceptance & 本次 run 满足其预登记的具体验收标准；未自动满足独立研究条件。\\
\code{operational_ready} & 运行产物和请求记录足够完整且专项验收没有运行阻断。\\
\code{claim_ready} & 在 operational ready 之上，还需合格 final 角色、显式 claim eligible、干净记账和独立评价前提。\\
\bottomrule
\end{tabularx}
\caption{主要验收关口。不同层级分别记录与报告。}
\end{table}

有 trusted fixture 替代独立评价、数据角色非 final、usage 冲突或缺失、解析错误等情况会影响对应关口。当前开发门禁即使全部通过，也不会把 \code{claim_eligible=false} 自动改为 true。
\loc{evaluation/protocols/readiness.py; harness/validation.py}

\subsection{评价指标与成本定义}
三值结果保留弃权及失败分母，成本分别按上下文、模型和执行报告。
\subsubsection{保留完整分母}
设全部评价项数为 $N$，实际正例与负例数为 $N_+,N_-$。实质判断的混淆计数分别为 $TP$、$FP$、$TN$ 和 $FN$。弃权、失败、中断和未运行作为独立计数保留，而不伪装成正常 SAFE 或从总体分母删除。当前主要指标为
\begin{align}
\mathrm{Coverage}&=\frac{TP+FP+TN+FN}{N},\\
\mathrm{StrictRecall}&=\frac{TP}{N_+},\qquad
\mathrm{CoveredRecall}=\frac{TP}{TP+FN},\\
\mathrm{Precision}&=\frac{TP}{TP+FP},\\
\mathrm{PopulationFPR}&=\frac{FP}{N_-},\qquad
\mathrm{CoveredFPR}=\frac{FP}{FP+TN},\\
\mathrm{ConservativeFPR}&=\frac{N_--TN}{N_-},\\
\mathrm{Accuracy}_{\mathrm{all}}&=\frac{TP+TN}{N}.
\end{align}
分母为零时返回未定义值。严格召回和全量 accuracy 保留未覆盖项影响；covered 指标回答已经形成实质结论的子集表现。保守 FPR 将所有未证成真阴性的负例计入保守记分，应与真实观测 FP 分开解释。

部分源码发现计分只具有正例参考，因此 FPR 明确为未定义。配对的真假标签判断、关键位置发现命中、advisory 召回、具备匹配 validator 的票数，以及实际执行见证次数，都有不同分母，不能互相替换。
\loc{evaluation/metrics.py; evaluation/discovery.py; harness/results.py}

\subsubsection{区间与配对比较}
metrics 模块提供描述性的 Wilson 二项区间，并明确它不是跨仓库泛化区间。缺失项、中断或未完成的 provisional 结果保留标识，不能通过只挑完整格子来制造确定性结论。

配对统计按 case 与 system 对齐，记录左赢、右赢、相同、不可用及标签分歧。案例和同仓库多个入口具有相关性；单轮开发对照没有自动获得重复实验或仓库级独立样本的解释权。显著成本或召回改进需要另行冻结重复次数、统计单位和分析规则。

\subsubsection{成本的三个维度}
\textbf{上下文成本}记录检索和观察的序列化文本上界；\textbf{模型成本}记录实际请求与 provider reported usage；\textbf{执行成本}记录工具尝试、物理 witness、资源和墙钟。三者在结果中共存，而不是把一个工具观察单位换算成一个 provider token。

E5 固定票回放可计算哪些后续任务会被跳过，但该值是反事实节省，未测量新调度下模型回复、工具耗时或总 wall-time。Jev 也需要将 routing 请求、source read 和 scoped planner 成本计入总额；最终标签正确不单独证明 routing 已采用或节省了请求。


\subsection{新的只读 journal 恢复审计}
\code{audit_journal} 对现存 JSONL 事件进行只读审计，区分完整事件、缺少换行的末行、损坏 JSON、未知或未完成请求、usage 冲突及尚无确定用量的 reservation。每个有问题的行记录行号和原始 bytes 摘要，避免将可疑内容直接回显。unknown usage 保留 unknown，不记为零；审计不退款、不补造响应、不自动重发。它给操作者一个保守的恢复清单，不提供事务式自动续跑保证。

新增控制还覆盖六个并发最后名额竞争者、畸形 model content 的 usage 保留、取消后不启动后续调用，以及已知 fake key 出现在 URL/HTTP 异常中的脱敏。错误类别、阶段和私有引用进入诊断；这不构成对任意未知秘密格式的完全识别。产品默认不加载 API 凭据，gate/model 阶段为零。\citep{acceptance}
\loc{runtime/journal.py::audit_journal; runtime/model.py; runtime/redaction.py}
